\section{Responsible Use of Large Language Models}

For this course, you are allowed to use \textbf{Large Language Models (LLMs)}, in fact recommended to use. LLMs can support brainstorming, drafting, summarising background information, or proposing alternative ways to express an idea. They may be useful when your team is stuck, when you need help converting rough notes into a clearer structure, or when you want to compare several possible approaches quickly.

However, LLM output must never replace your own engineering judgement. In this module, you remain responsible for the accuracy, relevance, and originality of everything you submit. When using an LLM, your team should:
\begin{itemize}
    \item disclose that the tool was used and describe how it supported the work;
    \item verify technical claims, calculations, and references against credible sources;
    \item rewrite and adapt the output so that it reflects your own understanding and project context; and
    \item keep evidence of important prompts and outputs when the lecturer asks how the work was produced.
\end{itemize}

For customer-needs work, direct interaction with real users remains the preferred method because it produces richer evidence than a simulated answer. If an LLM is used to generate starting points, treat those outputs as drafts that must be validated through observation, interviews, benchmarking, or literature.

Careful and transparent LLM use can save time. Uncritical copying, however, usually leads to weak reports because the output often ignores the exact question, the technical constraints of the product, or the level of evidence expected in an Industrial Physics course.
